\documentclass[journal]{IEEEtran}

% ---- Packages ----
\usepackage[utf8]{inputenc}
\usepackage[T1]{fontenc}
\usepackage{amsmath,amssymb,amsfonts,amsthm}
\usepackage{graphicx}
\usepackage{textcomp}
\usepackage{xcolor}
\usepackage{booktabs}
\usepackage{hyperref}
\usepackage{url}
\usepackage{cite}

\hypersetup{
  colorlinks=true,
  linkcolor=black,
  citecolor=black,
  urlcolor=blue
}

% ---- Theorem environments ----
\newtheorem{theorem}{Theorem}
\newtheorem{lemma}[theorem]{Lemma}
\newtheorem{proposition}[theorem]{Proposition}
\newtheorem{corollary}[theorem]{Corollary}
\newtheorem{conjecture}[theorem]{Conjecture}
\theoremstyle{definition}
\newtheorem{definition}[theorem]{Definition}
\theoremstyle{remark}
\newtheorem{remark}[theorem]{Remark}
\newtheorem{example}[theorem]{Example}

% ---- Custom commands ----
\newcommand{\R}{\mathbb{R}}
\newcommand{\E}{\mathbb{E}}
\newcommand{\norm}[1]{\left\lVert #1 \right\rVert}
\newcommand{\ip}[2]{\langle #1, #2 \rangle}
\DeclareMathOperator*{\argmin}{arg\,min}
\DeclareMathOperator*{\argmax}{arg\,max}
\DeclareMathOperator{\tr}{tr}

\begin{document}

% ---- Title ----
\title{Geometric Aesthetics: A Framework for Beauty\\as
Compressibility in Perceptual Eigenspaces}

\author{Andrew~H.~Bond
\thanks{Department of Computer Engineering, San Jose State University.
E-mail: andrew.bond@sjsu.edu}}

\markboth{Geometric Aesthetics}{Bond: Beauty as Compressibility
in Perceptual Eigenspaces}

\maketitle

% =====================================================================
% ABSTRACT
% =====================================================================
\begin{abstract}
We propose a mathematical framework in which aesthetic
experience---the perception of beauty, harmony, and
elegance---arises from the compressibility of stimuli relative
to a perceiver's learned perceptual eigenstructure.  We formalize
perceptual spaces as finite-dimensional inner product spaces equipped
with a learned covariance, define a scoring functional that measures
description efficiency relative to this structure, and derive three
main results under simplifying assumptions:
(1)~the scoring functional is maximized at a characterizable
boundary between order and randomness;
(2)~stimuli with group symmetry achieve exact compression as a
special case;
(3)~a temporal extension yields a compression-progress functional
whose peaks correspond to moments of ``insight.''
We show how Birkhoff's aesthetic measure, Berlyne's inverted-U,
and reward prediction error theory can be interpreted within this
framework, and we derive heuristic predictions testable against
existing datasets.  The framework is proposed as a formalization of
intuitions, not as a closed theory; we are explicit about where
the mathematics is rigorous and where it is conjectural.
\end{abstract}

\begin{IEEEkeywords}
Aesthetics, information theory, perception, symmetry, compression,
eigenspectrum, beauty.
\end{IEEEkeywords}


% =====================================================================
\section{Introduction}
\label{sec:intro}
% =====================================================================

The question of why certain configurations are perceived as beautiful
has occupied philosophers since antiquity and scientists since
Fechner's experimental aesthetics in the 1870s~\cite{fechner1876vorschule}.
Despite sustained inquiry, no mathematical framework has achieved broad
acceptance.

We propose that aesthetic experience has a geometric structure that
can be partially formalized.  The central idea is simple: each
perceiver carries a \emph{perceptual eigenstructure}---a learned
decomposition of the stimulus space into directions of decreasing
importance---and the aesthetic response to a stimulus depends on the
relationship between the stimulus's structure and the perceiver's
eigenstructure.  A stimulus that is compressible in the perceiver's
eigenbasis but not in a random basis carries \emph{meaningful
structure relative to that perceiver}.  We identify this property
with beauty.

\paragraph{Scope and honesty.}
This paper is a \emph{conceptual theory paper with formal
scaffolding}, not a theorem-driven paper in the classical sense.
We prove some results rigorously
(Propositions~\ref{prop:expertise}, \ref{prop:symmetry_compression},
\ref{prop:existence}, and~\ref{prop:berlyne_shape}).
Other results are heuristic observations or conjectures, labeled
as such (the temporal theory in~\S\ref{sec:temporal}, broken
symmetry in~\S\ref{sec:symmetry}, connections to neuroscience
and cultural distance in~\S\ref{sec:conjectural}).  We believe the framework is
productive even where it is not yet rigorous, but we do not claim
more rigor than exists.


% =====================================================================
\section{The Perceptual Eigenspace}
\label{sec:eigenspace}
% =====================================================================

\subsection{Setup}

Let $\mathcal{X} \subseteq \R^d$ be a \emph{stimulus space}---the
set of all possible stimuli in a domain.  A perceiver~$P$ has been
exposed to a distribution~$\mu$ over~$\mathcal{X}$ (their lifetime
experience in the domain).

\begin{definition}[Perceptual covariance]
\label{def:cov}
The \emph{perceptual covariance} of perceiver~$P$ is the positive
semi-definite matrix
\begin{equation}
    \Sigma_P = \E_\mu\bigl[(x - \bar{x})(x - \bar{x})^\top\bigr],
    \quad \bar{x} = \E_\mu[x].
    \label{eq:cov}
\end{equation}
Its eigendecomposition $\Sigma_P = U \Lambda U^\top$ with
$\lambda_1 \geq \cdots \geq \lambda_d \geq 0$ defines the
perceiver's \emph{perceptual eigenbasis}.
\end{definition}

\begin{definition}[Perceptual coordinates]
\label{def:coords}
The representation of stimulus~$x$ in the perceiver's eigenbasis is
$z = U^\top(x - \bar{x}) \in \R^d$.  In these coordinates,
$\Sigma_P$ is diagonal with entries~$\lambda_i$.
\end{definition}

\begin{definition}[Participation ratio]
\label{def:pr}
For a nonnegative spectrum $s_1, \ldots, s_d$ with $\sum s_i > 0$,
the \emph{participation ratio} is
\begin{equation}
    D_{\mathrm{eff}}(s) = \frac{(\sum_{i=1}^d s_i)^2}
    {\sum_{i=1}^d s_i^2},
    \label{eq:pr}
\end{equation}
satisfying $1 \leq D_{\mathrm{eff}} \leq d$.  Applied to the
eigenvalues, $D_{\mathrm{eff}}(\lambda)$ measures the effective
dimensionality of the perceiver's experience.  Applied to
$\{z_i^2\}$, it measures the effective dimensionality of a stimulus.
\end{definition}

\begin{proposition}[Expertise can expand effective dimensionality]
\label{prop:expertise}
Let $\Sigma' = (1-\alpha)\Sigma + \alpha \Sigma_{\mathrm{new}}$
for $\alpha \in (0,1)$.  If the normalized eigenvalue spectrum of
$\Sigma'$ majorizes that of $\Sigma$ less sharply---that is, if the
update \emph{flattens} the normalized spectrum in the sense that
$\sum_{i=1}^k \lambda'_i / S' \leq \sum_{i=1}^k \lambda_i / S$
for all $k < d$ (where $S = \sum \lambda_i$, $S' = \sum \lambda'_i$,
and eigenvalues are sorted descending)---then
$D_{\mathrm{eff}}(\Sigma') \geq D_{\mathrm{eff}}(\Sigma)$,
with equality iff the normalized spectra coincide.
\end{proposition}

\begin{proof}
The participation ratio is $D_{\mathrm{eff}} = 1 / \sum q_i^2$
where $q_i = \lambda_i / S$ is the normalized spectrum.  The map
$q \mapsto \sum q_i^2$ is a Schur-convex function on the probability
simplex.  By the Schur--Ostrowski criterion, $\sum (q'_i)^2 \leq
\sum q_i^2$ whenever $q'$ is majorized by $q$ (i.e., the
normalized spectrum of $\Sigma'$ is more uniform than that of
$\Sigma$).  Therefore $D_{\mathrm{eff}}(\Sigma') \geq
D_{\mathrm{eff}}(\Sigma)$.  Equality holds iff $q = q'$.
\end{proof}

\begin{remark}
The majorization condition is not automatically satisfied by
adding variance on low-eigenvalue directions; it depends on the
specific update.  The proposition says that expertise
\emph{can} expand effective dimensionality---and does so precisely
when the learning flattens the normalized spectrum.  This is a
sufficient condition, not a claim that all learning increases
$D_{\mathrm{eff}}$.
\end{remark}

This formalizes the intuition that expertise in a domain expands
the perceiver's effective perceptual dimensionality: a wine expert
discriminates along dimensions invisible to a novice.


% =====================================================================
\section{The Aesthetic Score}
\label{sec:score}
% =====================================================================

\subsection{Motivation and Design}

We seek a score $\mathcal{A}(x; \Sigma)$ that captures: \emph{a
stimulus is aesthetically valued when it carries meaningful
structure relative to the perceiver}.  ``Meaningful structure'' means
the stimulus is (i)~not trivial (it has content across multiple
dimensions), (ii)~not random (it is compressible), and
(iii)~compressible specifically in the perceiver's eigenbasis
(the structure is aligned with learned experience).

\subsection{Definition}

We work on the unit energy surface $\norm{z}^2 = 1$ to eliminate
scale dependence.  All stimuli are normalized to unit energy in
perceptual coordinates before scoring.  We restrict to the
\emph{active subspace}: let $r = |\{i : \lambda_i > 0\}|$ be the
rank of $\Sigma_P$, and work in the $r$-dimensional subspace
spanned by eigenvectors with positive eigenvalues.  This avoids
division by zero in the compressibility term.

\begin{definition}[Aesthetic score]
\label{def:score}
Let $z \in \R^r$ with $\norm{z}^2 = 1$, where $r$ is the rank of
$\Sigma_P$ and all coordinates correspond to eigenvectors with
$\lambda_i > 0$.  Define the energy distribution $p_i = z_i^2$
(so $p_i \geq 0$ and $\sum_{i=1}^r p_i = 1$).  Let
$q_i = \lambda_i / \sum_{j=1}^r \lambda_j$ be the normalized
perceiver eigenvalues.  The \emph{aesthetic score} is
\begin{equation}
    \mathcal{A}(p; q) =
    \underbrace{D_{\mathrm{eff}}(p)}_{\text{complexity}}
    \;\cdot\;
    \underbrace{\exp\!\left(-\frac{1}{r}\sum_{i=1}^r
    \frac{p_i}{q_i}\right)}_{\text{compressibility}},
    \label{eq:score}
\end{equation}
which is well-defined since $q_i > 0$ for all $i \leq r$.
\end{definition}

The first factor, $D_{\mathrm{eff}}(p) = 1/\sum p_i^2$, measures
how many dimensions the stimulus occupies (complexity).  The second
factor measures how well the stimulus's energy distribution aligns
with the perceiver's eigenvalue distribution: when $p_i$ concentrates
on high-$\lambda_i$ directions, the argument of the exponential is
small and the compressibility factor is close to~1; when energy is
spread onto low-$\lambda_i$ directions, the factor decreases.

The exponential form ensures the compressibility factor is bounded
in $[0, 1]$, eliminating the sign and scale issues of the linear
penalty used in the preliminary version.

\begin{remark}
On the unit simplex $\{p : p_i \geq 0, \sum p_i = 1\}$, the
score $\mathcal{A}(p; \lambda)$ is well-defined, continuous,
and bounded.  At $p = e_i$ (all energy on one dimension),
$D_{\mathrm{eff}} = 1$ and $\mathcal{A}$ is small.  At
$p = (1/d, \ldots, 1/d)$ (uniform), $D_{\mathrm{eff}} = d$
but the compressibility factor is generally small for non-uniform
$\lambda$.  The maximum occurs at an interior point of the simplex.
\end{remark}


% =====================================================================
\section{The Aesthetic Optimum}
\label{sec:optimum}
% =====================================================================

\begin{proposition}[Existence of an interior maximizer]
\label{prop:existence}
The aesthetic score $\mathcal{A}(p; \lambda)$ on the closed simplex
$\Delta^{d-1} = \{p : p_i \geq 0, \sum p_i = 1\}$ attains its
maximum.  The maximizer does not lie at any vertex $p = e_i$
(provided $d \geq 2$ and at least two $\lambda_i > 0$).
\end{proposition}

\begin{proof}
$\mathcal{A}$ is continuous on the compact set $\Delta^{d-1}$, so
attains its maximum by the extreme value theorem.  At any vertex
$p = e_i$, $D_{\mathrm{eff}}(e_i) = 1$, so $\mathcal{A}(e_i;
\lambda) = 1 \cdot \exp(-1/(d \cdot \lambda_i / S))$ where
$S = \sum \lambda_j$.  By moving a small amount of mass $\delta$
from coordinate $i$ to a coordinate $j$ with $\lambda_j > 0$,
$D_{\mathrm{eff}}$ increases to first order (since
$\partial D_{\mathrm{eff}}/\partial p_j |_{p=e_i}$ is unbounded
from above as we leave the vertex), while the exponential factor
changes continuously.  Therefore $e_i$ is not a maximizer.
\end{proof}

\begin{proposition}[Inverted-U on the power-law family]
\label{prop:berlyne_shape}
On the one-parameter family $p_i(\alpha) = c_\alpha \cdot i^{-\alpha}$
(where $c_\alpha$ normalizes to unit sum), the score
$\mathcal{A}(\alpha) = \mathcal{A}(p(\alpha); \lambda)$ satisfies:
\begin{enumerate}
    \item[(a)] $\mathcal{A}(0) > 0$ and $\mathcal{A}(\alpha) \to 0$
    as $\alpha \to \infty$.
    \item[(b)] There exists $\alpha^* > 0$ achieving the maximum.
    \item[(c)] As $\beta \to 0$ (flat perceiver spectrum),
    $\alpha^* \to 0$.
\end{enumerate}
That is, the score is an inverted-U function of concentration
(and hence of stimulus complexity).
\end{proposition}

\begin{proof}
(a) At $\alpha = 0$, $p_i = 1/d$ and $D_{\mathrm{eff}} = d$, so
$\mathcal{A}(0) = d \cdot \exp(-S^{-1} \sum \lambda_i^{-1} / d^2)
> 0$.  As $\alpha \to \infty$, all mass concentrates on $i = 1$, so
$D_{\mathrm{eff}} \to 1$ and $\mathcal{A} \to \exp(-1/(d \cdot
\lambda_1 / S))$, which is less than $\mathcal{A}(0)$ for $d \geq 2$.

(b) By continuity of $\mathcal{A}(\alpha)$ on $[0, \infty)$ and
the boundary behavior in~(a), the maximum is attained at some
$\alpha^* \in (0, \infty)$.

(c) When $\beta = 0$ (all $\lambda_i$ equal), the compressibility
factor is constant for all $p$, so $\mathcal{A}$ is proportional to
$D_{\mathrm{eff}}(p)$, which is maximized at $\alpha = 0$.
\end{proof}

This formalizes Berlyne's~\cite{berlyne1971aesthetics} central
finding: the optimal stimulus complexity is intermediate.

\begin{remark}[What we do not prove]
We do not prove that $\alpha^* < \beta$ in general.  Numerical
experiments suggest this holds for moderate $\beta$ (the aesthetically
optimal stimulus decays more slowly than the perceiver's spectrum),
but for sharply peaked spectra ($\beta \gtrsim 2$ at moderate $d$)
the maximizer can occur at $\alpha > \beta$.  We also do not prove
uniqueness of $\alpha^*$ or characterize the global optimum over the
full simplex.  These remain open problems.
\end{remark}


% =====================================================================
\section{Symmetry as Compression}
\label{sec:symmetry}
% =====================================================================

The classical theory of beauty as symmetry is a special case of
compression.

\begin{definition}[Fixed-point subspace]
Let $G$ be a finite group acting on $\R^d$ via linear
transformations $\{T_g\}_{g \in G}$.  The \emph{fixed-point
subspace} is $\mathrm{Fix}(G) = \{v \in \R^d : T_g v = v
\;\forall g \in G\}$.
\end{definition}

\begin{proposition}[Symmetry implies compression]
\label{prop:symmetry_compression}
If $x \in \mathrm{Fix}(G)$ for some representation of $G$ on
$\R^d$, then $x$ is determined by $\dim(\mathrm{Fix}(G))$
coordinates.  The compression ratio relative to the ambient space
is $d / \dim(\mathrm{Fix}(G))$.
\end{proposition}

\begin{proof}
$\mathrm{Fix}(G)$ is a linear subspace of $\R^d$.  Any element
$x \in \mathrm{Fix}(G)$ can be represented in a basis for this
subspace using $\dim(\mathrm{Fix}(G))$ coordinates.  The compression
ratio follows by definition.
\end{proof}

\begin{remark}
The dimension $\dim(\mathrm{Fix}(G))$ depends on the specific
representation, not just on $|G|$.  For the regular representation,
$\dim(\mathrm{Fix}(G)) = 1$ (maximal compression).  For a
representation with many trivial summands, $\dim(\mathrm{Fix}(G))$
can be much larger.  Specific examples:
\end{remark}

\begin{example}[Bilateral symmetry]
For $G = \mathbb{Z}/2\mathbb{Z}$ acting on $\R^d$ by
$(z_1, z_2, \ldots, z_{d/2}, z_{d/2+1}, \ldots, z_d) \mapsto
(z_{d/2+1}, \ldots, z_d, z_1, \ldots, z_{d/2})$ (swapping left and
right halves), $\mathrm{Fix}(G)$ consists of vectors with identical
halves: $\dim(\mathrm{Fix}(G)) = d/2$, giving $2:1$ compression.
\end{example}

\begin{example}[Rotational symmetry of order $n$]
For $G = \mathbb{Z}/n\mathbb{Z}$ acting by cyclic permutation on
$\R^n$, $\mathrm{Fix}(G)$ consists of constant vectors:
$\dim(\mathrm{Fix}(G)) = 1$, giving $n:1$ compression.
\end{example}

\paragraph{Broken symmetry (heuristic).}
The intuition that perfect symmetry is boring---and that small
asymmetries can enhance aesthetic value---is well-supported
empirically (slightly asymmetric faces are preferred over perfectly
symmetric ones).  Within our framework, this corresponds to the
observation that a small perturbation $\epsilon \perp \mathrm{Fix}(G)$
added to a symmetric stimulus $x_s \in \mathrm{Fix}(G)$ can increase
$D_{\mathrm{eff}}$ (by activating new dimensions) faster than it
decreases the compressibility factor (by spreading energy onto
less-aligned directions).

However, whether $\mathcal{A}(x_s + \epsilon) > \mathcal{A}(x_s)$
depends on the specific directions in which $\epsilon$ lies.  If
the perturbation projects onto directions with very small $\lambda_i$
(directions the perceiver considers unimportant), the compressibility
penalty can dominate.  A general proposition would require an
explicit condition on the alignment of $\epsilon$ with the
perceiver's eigenspectrum.  We leave this as an open problem and
note that the phenomenon is robustly observed in the empirical
literature on symmetry preference~\cite{ramachandran1999science}.


% =====================================================================
\section{Compression Progress: A Temporal Extension}
\label{sec:temporal}
% =====================================================================

Many aesthetic experiences are temporal: music, narrative,
mathematical proof.  We propose a temporal extension, noting
that this section is more speculative than the preceding ones.

\subsection{Setup}

Let $x_1, x_2, \ldots, x_T \in \R^d$ be a sequence of stimuli.
We represent the perceiver's cumulative experience at time $t$ by the
sufficient statistic
\begin{equation}
    \bar{z}_t = \frac{1}{t}\sum_{s=1}^t z_s, \qquad
    S_t = \frac{1}{t}\sum_{s=1}^t z_s z_s^\top,
    \label{eq:sufficient}
\end{equation}
where $z_s = U^\top(x_s - \bar{x})$ are perceptual coordinates.
The sample covariance is $\hat{\Sigma}_t = S_t - \bar{z}_t
\bar{z}_t^\top$.

\begin{definition}[Compression progress]
\label{def:progress}
Let $p_s^{(t)}$ denote the normalized energy distribution of $z_s$
relative to the eigenbasis of $\hat{\Sigma}_t$: that is,
$p_{s,i}^{(t)} = (u_i^{(t)\top} z_s)^2 / \norm{z_s}^2$
where $\{u_i^{(t)}\}$ are the eigenvectors of $\hat{\Sigma}_t$,
and let $q^{(t)}$ be the normalized eigenvalues of $\hat{\Sigma}_t$.
The \emph{compression progress} at time $t$ is
\begin{equation}
    \dot{\mathcal{A}}_t =
    \sum_{s=1}^t \mathcal{A}(p_s^{(t)}; q^{(t)})
    - \sum_{s=1}^t \mathcal{A}(p_s^{(t-1)}; q^{(t-1)}).
    \label{eq:progress}
\end{equation}
This measures how much better the perceiver can describe all prior
stimuli using the updated model versus the previous model.
\end{definition}

In words: compression progress measures how much better the
perceiver can describe \emph{all data seen so far} using their
updated model.  A large positive value at time $t$ means $x_t$
caused a model update that retroactively made everything more
compressible---a moment of insight.

\begin{conjecture}[Compression progress under rank-one updates]
\label{conj:progress}
If $\hat{\Sigma}_t$ differs from $\hat{\Sigma}_{t-1}$ by an
approximate rank-one update (the new stimulus is the dominant
source of change), then $\dot{\mathcal{A}}_t$ is approximately
proportional to the squared projection of $z_t$ onto the
eigenvector of $\hat{\Sigma}_{t-1}$ that undergoes the largest
relative eigenvalue increase.

In particular, compression progress is large when the new stimulus
(i)~projects strongly onto a previously low-variance direction, and
(ii)~reveals a systematic pattern (not just noise) that makes prior
data more compressible along that direction.
\end{conjecture}

We do not prove this conjecture in full generality.  The difficulty
is that the aesthetic score $\mathcal{A}$ involves the participation
ratio, which is not differentiable at certain boundaries and does not
decompose cleanly under rank-one updates.  A rigorous treatment would
require smoothing or restricting to the power-law family.  We include
the conjecture because it generates testable predictions and
connects to the neuroscience literature.

\begin{example}[The punchline of a joke]
A joke's setup builds a model (high-$\lambda$ directions).  The
punchline reveals an alternative interpretation: a previously
zero-variance direction suddenly explains the setup data.  The
compression progress is large because the entire setup sequence
becomes more compressible under the new model.
\end{example}

\begin{example}[Modulation in music]
A modulation from the home key to a distant key projects onto a
previously low-variance direction in harmonic space.  The return
to the home key causes another spike: the listener's model updates
to accommodate the overall tonal plan.
\end{example}


% =====================================================================
\section{Connections to Classical Theories}
\label{sec:connections}
% =====================================================================

We now show how the framework \emph{relates to} (not ``subsumes'')
three classical theories.

\paragraph{Birkhoff's aesthetic measure.}
Birkhoff~\cite{birkhoff1933aesthetic} defined $M = O/C$ (order over
complexity).  In our framework, order corresponds to compression via
symmetry (Proposition~\ref{prop:symmetry_compression}) and complexity
to $D_{\mathrm{eff}}$.  The ratio $O/C$ is a monotone function of
the compressibility-per-dimension, which is one component of our
score.  Our framework replaces the division (which penalizes
complexity) with a product (which rewards complexity up to a point),
addressing Birkhoff's known failure to account for the appeal of
complex stimuli.

\paragraph{Berlyne's inverted-U.}
Corollary~\ref{prop:berlyne_shape} derives the inverted-U as a
consequence of the competing drives toward complexity and
compressibility.  The novel prediction is that the peak shifts with
perceiver eigenvalue distribution, which is measurable through
preference experiments stratified by expertise.

\paragraph{Reward prediction error.}
The neuroscience literature reports that aesthetic pleasure
correlates with dopaminergic prediction
error~\cite{salimpoor2011anatomically}.  We conjecture
(but do not prove) that compression progress
(Definition~\ref{def:progress}) is a computational-level description
of the same signal: the brain's reward system fires when the
internal model improves---when something that was complex becomes
compressible.  This interpretation follows
Schmidhuber~\cite{schmidhuber2009simple}, who first proposed
compression progress as a theory of beauty.  Our contribution is to
ground it in a specific eigenspace framework rather than abstract
Kolmogorov complexity.


% =====================================================================
\section{Conjectural Extensions}
\label{sec:conjectural}
% =====================================================================

We briefly note two extensions that are speculative but
suggestive.

\paragraph{Cultural distance.}
The space of positive definite matrices $\mathcal{S}_{++}^d$
carries the Fisher information metric, under which the geodesic
distance between $\Sigma_A$ and $\Sigma_B$ is
$d_F = \norm{\log(\Sigma_A^{-1/2}\Sigma_B\Sigma_A^{-1/2})}_F$.
We conjecture that aesthetic disagreement between cultural groups
$A$ and~$B$ on a stimulus~$x$ is bounded by a Lipschitz multiple
of $d_F(\Sigma_A, \Sigma_B)$.  This would formalize the intuition
that cultures with similar artistic traditions have similar aesthetic
preferences.  A proof would require careful handling of the
coordinate-dependent nature of the score.

\paragraph{Cross-domain universality.}
If the same perceiver has eigenstructures in different
domains (visual, musical, gustatory), the framework predicts that
aesthetic principles are \emph{structurally identical} across
domains but the \emph{eigenspaces are different}.  Symmetry is
beautiful in faces and in music for the same reason (compression),
but a visual art expert's expanded eigenbasis does not transfer to
music appreciation.  This is testable.


% =====================================================================
\section{Predictions}
\label{sec:predictions}
% =====================================================================

The framework generates heuristic predictions, all testable against
existing or obtainable data:

\begin{enumerate}
    \item \textbf{Expertise shift.}  Experts prefer stimuli with
    higher effective dimensionality than novices
    (Corollary~\ref{prop:berlyne_shape}).  Testable by comparing
    preferred complexity in trained vs untrained listeners.

    \item \textbf{Broken symmetry preference.}  Stimuli with
    small controlled asymmetry should be preferred over both
    perfect symmetry and strong asymmetry, provided the
    perturbation aligns with perceiver-important directions.
    Testable against face preference data with parametric
    asymmetry manipulation.

    \item \textbf{Compression progress predicts chills.}
    Musical ``frisson'' should correlate with peaks of
    $\dot{\mathcal{A}}_t$ (Conjecture~\ref{conj:progress}).
    Testable by computing the score from MIDI and comparing
    against physiological chill recordings.

    \item \textbf{Mathematical elegance.}  Among proofs of the
    same theorem, mathematicians should rate shorter proofs (lower
    description length) as more elegant.  Approximately testable
    using proof length in a fixed formal system.

    \item \textbf{No cross-domain transfer.}  Training in visual
    art should not shift preferred complexity in music, because the
    eigenspaces are domain-specific.  Testable via a pre/post
    design with domain-specific preference measurement.
\end{enumerate}


% =====================================================================
\section{Discussion}
\label{sec:discussion}
% =====================================================================

\paragraph{What is rigorous.}
Propositions~\ref{prop:expertise} (under the majorization condition),
\ref{prop:symmetry_compression}, \ref{prop:existence}, and
\ref{prop:berlyne_shape} are proved.  The existence of an interior
maximizer and the inverted-U shape on the power-law family are
established; the location and uniqueness of the maximizer are not.

\paragraph{What is heuristic.}
The aesthetic score (Definition~\ref{def:score}) is a motivated
construction, not derived from first principles.  The exponential
compressibility factor is chosen for boundedness and mathematical
convenience; other monotone, bounded functions of the alignment
penalty would yield qualitatively similar behavior.  The broken
symmetry phenomenon (\S\ref{sec:symmetry}) is a heuristic
observation supported by empirical evidence but not proved in
full generality.

\paragraph{What is conjectural.}
The temporal theory (Conjecture~\ref{conj:progress}), the cultural
distance bound, the connection to reward prediction error, and the
cross-domain universality claim are interpretive connections, not
proved results.  We include them because they generate testable
predictions and point toward a research program.

\paragraph{Linearity.}
We model perceptual spaces as linear.  Real perceptual spaces are
likely nonlinear manifolds.  Extending to kernel PCA or
autoencoders would capture nonlinear structure at the cost of
analytical tractability.  The qualitative predictions
(inverted-U, expertise shift, broken symmetry) should survive
nonlinearization, but the closed-form results will not.

\paragraph{The PCA-Matryoshka connection.}
In PCA-Matryoshka~\cite{bond2026pcamatryoshka}, we showed that
eigenvalues serve as importance weights for embedding
compression: truncating in the PCA basis preserves meaning while
truncating in a random basis destroys it.  The aesthetic framework
developed here suggests this may reflect something deeper about
perception.  Beauty may be the subjective experience of encountering
meaningful structure---structure that compresses in the basis
evolution and experience gave us to understand the world.

This connection is metaphorical, not proved.  But we find it
suggestive enough to state.


% =====================================================================
\section{Conclusion}
\label{sec:conclusion}
% =====================================================================

We have proposed a mathematical framework for aesthetic experience
based on compressibility in perceptual eigenspaces.  The framework
yields four proved results (expertise expands effective
dimensionality under a majorization condition; symmetry compresses
description length; an interior maximizer exists; the score is
inverted-U on the power-law family), one heuristic observation
(broken symmetry), one conjecture (compression progress for temporal
aesthetics), and five testable predictions.

The framework interprets---but does not subsume---Birkhoff's
aesthetic measure, Berlyne's inverted-U, and reward prediction
error theory.  The main value, we believe, is not in the specific
functional form chosen, but in the conceptual identification:
\emph{the perceiver's learned eigenstructure determines what counts
as meaningful structure, and beauty is the perception of structure
that is meaningful in this sense}.

The framework is incomplete.  A complete theory of aesthetics
would need to account for temporal dynamics beyond our conjecture,
social and cultural mechanisms beyond our Fisher metric speculation,
emotional valence beyond our compression-based proxy, and the
irreducible role of personal history and association.  We offer this
framework not as a final word but as a mathematical starting point
that is precise enough to be wrong in testable ways.


% ---- References ----
\bibliographystyle{IEEEtran}
\begin{thebibliography}{15}

\bibitem{bond2026pcamatryoshka}
A.~H.~Bond, ``PCA-Matryoshka: Enabling effective dimension reduction
for non-Matryoshka embedding models,'' \emph{IEEE Trans.\ Artificial
Intelligence}, 2026.

\bibitem{birkhoff1933aesthetic}
G.~D.~Birkhoff, \emph{Aesthetic Measure}.\hskip 1em plus 0.5em minus
0.4em Cambridge, MA: Harvard Univ.\ Press, 1933.

\bibitem{berlyne1971aesthetics}
D.~E.~Berlyne, \emph{Aesthetics and Psychobiology}.\hskip 1em plus
0.5em minus 0.4em New York: Appleton-Century-Crofts, 1971.

\bibitem{schmidhuber2009simple}
J.~Schmidhuber, ``Simple algorithmic theory of subjective beauty,
novelty, surprise, interestingness, attention, curiosity, creativity,
art, science, music, jokes,'' \emph{J.\ Artif.\ General Intelligence},
vol.~1, no.~1, pp.~1--32, 2009.

\bibitem{ramachandran1999science}
V.~S.~Ramachandran and W.~Hirstein, ``The science of art,''
\emph{J.\ Consciousness Studies}, vol.~6, no.~6--7, pp.~15--51, 1999.

\bibitem{tymoczko2011geometry}
D.~Tymoczko, \emph{A Geometry of Music}.\hskip 1em plus 0.5em minus
0.4em Oxford, UK: Oxford Univ.\ Press, 2011.

\bibitem{chatterjee2014neuroaesthetics}
A.~Chatterjee and O.~Vartanian, ``Neuroaesthetics,'' \emph{Trends
Cogn.\ Sci.}, vol.~18, no.~7, pp.~370--375, 2014.

\bibitem{salimpoor2011anatomically}
V.~N.~Salimpoor \emph{et al.}, ``Anatomically distinct dopamine release
during anticipation and experience of peak emotion to music,''
\emph{Nature Neurosci.}, vol.~14, no.~2, pp.~257--262, 2011.

\bibitem{fechner1876vorschule}
G.~T.~Fechner, \emph{Vorschule der Aesthetik}.\hskip 1em plus 0.5em
minus 0.4em Leipzig: Breitkopf \& H\"artel, 1876.

\bibitem{weyl1952symmetry}
H.~Weyl, \emph{Symmetry}.\hskip 1em plus 0.5em minus 0.4em Princeton,
NJ: Princeton Univ.\ Press, 1952.

\end{thebibliography}

\end{document}
